\subsection{Why quantum computing?}

Interest in quantum computing is high, reaching well beyond academic interest,
with many extant public and private sector projects seeking to develop, display,
and eventually exploit quantum computation. This interest motivates a great deal
of intellectual and economic investment in physics, computation, and information
science. On the commercial side IBM, NTT, Intel, and Toshiba lead the way in
patent applications as of 2020\cite{CommercialLandscape}. Additionally,
governments around the world have invested the equivalent of billions of USD in
initiatives on quantum technologies, such as the European Union, France,
Germany, Japan, Canada, and the US among others\cite{DownToBusiness}.

Attempting to prophesy on the future feasibility of useful quantum computation
is well beyond the scope of this document and the author's knowledge and
abilities, of course. John Preskill, a leading researcher in this field, labels
the current era of quantum computing to be the NISQ era, the era of \emph{noisy
intermediate-scale quantum} computing\cite{PreskillNISQ}. As such we are not
able yet able to leverage quantum computing to accomplish many of the most
referenced applications of quantum computing.

However, the author feels that such global interest in the fundamental ideas of
quantum mechanics and in understanding of computation and information beyond the
classical represents an exciting chance to learn ever more about the deeper
nature of our universe and reality. 

Our understanding of the world is not well posed to process quantum information.
To some degree, this can be reflected in our idea of "measurement" within
quantum mechanics. Consider the Stern-Gerlach experiment\cite{SternGerlach}.
After a stream of silver atoms are passed through an appropriately directed
inhomogenous magnetic field, the path of the silver atoms will be deflected
according to relative orientation between the magnetic field and the atom's
magnetic moment. The path of the atom may be recorded as a spot on the receiving
plate. When the experiment is performed many times two clear clusters of spots
emerge, corresponding respectively to the orientation of the atom's magnetic
moment completely along or against the external magnetic field. This was an
early success of quantum mechanics, as the theory predicted such an outcome
would be seen. Both quantum mechanics and classical mechanics predict that,
prior to the interaction of the silver atom with the external magnetic field,
the atom exists in a state of arbitrary orientation. Classically, an atom in
such a state would have a component of its magnetic moment lying along or
against the external magnetic field. The magnitude of this component is
arbitrary, and so is the degree of deflection. Quantum mechanics, on the other
hand, predicts the arbitrary initial state of the magnetic moment to take the
form of a superposed wavefunction with components lying entirely along or
against the external magnetic field. Such a two-level system is important to
quantum computing, as we will describe later, and is conventionally represented
as
\begin{equation}
    \ket{\state{\Psi}} = \alpha\ket{\state{0}} + \beta\ket{\state{1}}
\end{equation}
As such, the effect of the external magnetic field upon the moment of the atom
is discrete, as should be the position of the atom upon the detector plate. But
when does the superposition of states become one or the other? We can easily
formulate the two recorded positions of the atom on the detector plate as being
in a superposition. Does the resolution of the superposition happen when the
atom interacts with the magnetic field.  Perhaps it occurs in the human eye
recording the spot on the plate, or perhaps even further as result of visual
information processing in the human brain, or perhaps we have to resort to a
semi-mystical idea of a "consciousness" that measures the state of a human's
physical brain?

Such debates are well beyond the scope of this document, the author, and unless
clear experiments can be devised to test the implications of such
interpretations, they are firmly beyond the scope of physics. For our purposes,
we informally term "measurement" as the process by which classical information
is extracted from a quantum system. Perhaps anticipating the adage "shut up and
calculate!" variously attributed to Dirac or Feynman, Von Neumann discussed this
process and concluded that the exact place where the "measurement" takes place
has no impact upon the fundamental mathematics of the theory\cite{VonNeumann}.

Much of quantum computing's promise comes from the idea that quantum information
is a more robust concept then classical information, where an exchange of
information between particles at a point in time allows for more correlation
between those particles' states when later examined than is possible in a
classical setting, as discussed by Bell in his famous paper on Bell's
inequalities\cite{BellTheorem}. This property, known as \emph{entanglement}, can
be interpreted as an exchange of classical information in a nonlocal fashion, or
an example of reality's failure to be governed purely by classical information.
This indescribability of the quantum universe in classical terms may be asserted
in the form of a principle that might be called "counterfactual definiteness" or
in the idea associated with the so-called Everett Many-world hypothesis that
measurement fundamentally doesn't physically exist. A result to a classical
entity that has a single classical outcome has, in fact, many classical outcomes
that exist as part of the superposition of a larger quantum system including the
system at hand, its environment, and the observer. While the present author
subscribes somewhat to the "shut up and calculate!" adage, in the interests of
full disclosure, the present author has found the most conceptually useful
philosophical framework for anchoring his mathematical and physical work to be
Everett's many worlds with decoherence models for wave functions collapse. We
disclose this not to firmly assert a "correct belief" or factional affiliation
but because we want to aid communication with the reader. We anticipate that
while such philosophical sympathies might color our mode of description we
believe the results stand on their own physically and mathematically.

One of the ways in which quantum information might be said to be more robust
than classical information lies in the possibilities of encoding classical
information.  Firstly, quantum particles exist in a superposition of classically
ponderable states. We saw this above when we considered the state of the
Stern-Gerlach particle prior to measurement. To that end, we will formally
define the \emph{qubit}. Quantum error correction as a field recognizes a
difference between the simple physical qubit and the logical qubit. The logical
qubit is an assemblage of associated physical qubits working together to produce
a qubit of information that may be robustly involved in quantum computation
without unacceptable error. We focus in this document on the physical qubit. At
its core the physical qubit is a two-level system that may be described by a
two-term wavefunction of the form introduced above
\begin{equation}
\ket{\Psi}
    =
    \cos\left(\thetaup/2\right)\ket{0} + 
    \sin\left(\thetaup/2\right)\e^{\im\phiup}\ket{1}.
\end{equation}
With $\thetaup \in \left[0,\pi\right)$ and $\phiup \in \left(-\pi,+\pi\right]$
we can compare this to classical physical bits, which are likewise two-level
physical systems, but in the bit the system is either in the state $\ket{0}$ or
$\ket{1}$. We can already see that storing the precise state of a two-level
quantum system in a classical computer will require two complex numbers, or
about $256$ bits of storage. For the qubit and other two-level systems,
asserting a normalization condition and dropping an overall phase factor allows
us to represent a qubit using two real numbers $\thetaup$ and $\phiup$ as above.
Two real numbers are about $128$ bits in most computing systems. Already we can
see a clear difference between qubit and bit, but for one qubit it is not so
difficult to find $128$ spare bits to work with. We might be tempted to imagine
that a system of bits could be represented with two complex numbers for each
qubit. This is not the case. A linear increase in the number of qubits requires
an exponential increase in the number of bits due to the aforementioned
entanglement. The arbitrary states of a two-level quantum system are described
by a wavefunction of the form
\begin{equation}
    \ket{\Psi}
    =
    a_{00}\ket{00}
        +
    a_{01}\ket{01}
        +
    a_{10}\ket{10}
        +
    a_{11}\ket{11}
\end{equation}
which requires four complex numbers. Asserting normalization and arbitrary
overall phase only reduces the representation budget by $128$ bits, and becomes
increasingly irrelevant as system size increases. As we examine the
wavefunctions of larger and larger systems of qubits, we can see that the number
of complex numbers needed to represent the system goes as $\approx 2^{N}$. This
makes many-body systems with many quantum degrees of free very difficult to
simulate classically in general, though perhaps for individual problems
simplifications and analytic shortcuts can be found to reduce the computational
overhead. Still, bespoke solutions to computational expense just swaps computer
workhours for human ones, and generally human work hours are more expensive.

This then, is the most exciting prospect of quantum computing for the present
author. A truly digital, general quantum computer will make entire domains of
problems much more amenable to simulation and free some of those expensive human
workhours to survey much broader swathes of physics and physical understanding.

Beyond the promise of quantum simulation, the quantum computing project is
fundamentally a practical one, with much of its motivation focused on the
applications of a quantum computing device to specific problems. Some of the
well-known algorithms that are use to demonstrate a quantum advantage include
Shor's algorithm\cite{Shor} and Grover's search algorithm\cite{Grover}. Shor's
algorithm is a way to factor numbers, such as $15$ into $3$ and $5$, in
polynomial time. This is far from a trivial application -- easy factoring of
numbers has dramatic practical implications, particularly in RSA
cryptography\cite{RSA}. For RSA cryptography, a key component involves the idea
that the complexity for factoring a number scales expoenentially with the size
of the number. To date, no polynomial time algorithm to factor numbers for
classical computers has not been found. While it is has been proven that
classical polynomial factoring is impossible, it is not expected to be found
anytime in the near future. This is not the case for Shor's algorithm, which can
factor in polynomial time. As easy factoring would break RSA encryption, and as
RSA encryption underlies much of the secure communication over the internet, a
practical implementation of Shor's algorithm would have devasting security
conseqences as a zero day exploit.